% Revize 2026-10-27: cílené doplnění; nové úlohy a řešení jsou vlastní tvorba.
% Metodický vzor: 4_rocnik/pripravy/Uvod_do_kombinatoriky.pdf (samostatný dokument).
% Hodina 37, 1. ročník SPŠ. Rozsah: Rozpis_vyuky_MAT_1_rocnik.tex, řádek 37.

% Plán: 1_rocnik/tematicke_plany/SPS.pdf; společná projektová šablona.

% Učebnice: ucebnice/01_Zakladni_poznatky/Mocniny_a_odmocniny/Mocniny a Odmocniny.tex

% Původ úloh: Vlastní zadání a vlastní formulace výkladu; metodická inspirace zdroji.

% Metodický zdroj: externi_zdroje/realisticky/Mocniny/04_KISS.pdf

% URL: http://www.realisticky.cz/ucebnice/01%20Matematika%20S%C5%A0/01%20Z%C3%A1kladn%C3%AD%20poznatky/06%20Mocniny/04%20KISS.pdf

% Metodický zdroj: externi_zdroje/realisticky/Mocniny/09_Vypocty_s_mocninami.pdf

% URL: http://www.realisticky.cz/ucebnice/01%20Matematika%20S%C5%A0/01%20Z%C3%A1kladn%C3%AD%20poznatky/06%20Mocniny/09%20V%C3%BDpo%C4%8Dty%20s%20mocninami.pdf

% Metodický zdroj: externi_zdroje/realisticky/Odmocniny/07_Druha_odmocnina.pdf

% URL: http://www.realisticky.cz/ucebnice/01%20Matematika%20S%C5%A0/01%20Z%C3%A1kladn%C3%AD%20poznatky/02%20%C4%8C%C3%ADseln%C3%A9%20obory/07%20Druh%C3%A1%20odmocnina.pdf

\documentclass[11pt,a4paper]{article}

% Spolecny styl pro ucitelske pripravy MAT (Work/MAT)
% Verze 1.1, 2026-10-23
\NeedsTeXFormat{LaTeX2e}
\ProvidesPackage{priprava}[2026/10/23 v1.1 Sablona ucitelskych priprav MAT]

% Zakladni typografie a cestina
\RequirePackage[T1]{fontenc}
\RequirePackage[utf8]{inputenc}
\RequirePackage[czech]{babel}
\RequirePackage{lmodern}
\RequirePackage{microtype}
\RequirePackage{amsmath,amssymb}
\RequirePackage{xcolor}
\RequirePackage{enumitem}
\RequirePackage[most]{tcolorbox}
\RequirePackage{tabularx,booktabs,array}
\RequirePackage{fancyhdr}
\RequirePackage{needspace}
\RequirePackage{geometry}
\RequirePackage[hidelinks]{hyperref}

\geometry{a4paper,top=17mm,bottom=18mm,left=23mm,right=23mm}
\definecolor{notegray}{gray}{0.39}
\setlength{\parindent}{0pt}
\setlength{\parskip}{2.5pt plus 1pt minus 0.5pt}
\setlist[itemize]{topsep=2pt,itemsep=1pt,parsep=0pt,leftmargin=1.45em}
\setlist[enumerate]{topsep=2pt,itemsep=2pt,parsep=0pt,leftmargin=1.65em}
\setcounter{secnumdepth}{0}
\renewcommand{\arraystretch}{1.12}

% Zahlavi: vlevo rocnik | globalni tema, vpravo cislo: nazev hodiny
% Zapati zustava zcela prazdne.
\pagestyle{fancy}
\fancyhf{}
\renewcommand{\headrulewidth}{0.35pt}
\renewcommand{\footrulewidth}{0pt}
\setlength{\headheight}{14pt}
\newcommand{\nastavPripravu}[4]{%
  \fancyhead[L]{\small #1\enspace\textbar\enspace #2}%
  \fancyhead[R]{\small #3: #4}%
}

% Minimalisticke cernobile boxy. Kazda definice zvlast.
\tcbset{commonlesson/.style={enhanced,breakable=false,colback=white,colframe=black,
  colbacktitle=white,coltitle=black,boxrule=0.45pt,
  arc=0.8mm,outer arc=0.8mm,left=3mm,right=3mm,top=1.6mm,bottom=1.6mm,
  toptitle=1mm,bottomtitle=0.7mm,fonttitle=\bfseries,
  attach boxed title to top left={xshift=2.5mm,yshift=-1.15mm},
  boxed title style={colback=white,colframe=white,boxrule=0pt,left=1mm,right=1mm}}}
\newtcolorbox{zadaniBox}[1][]{commonlesson,title={Zadání},#1}
\newtcolorbox{prikladBox}[1][]{commonlesson,title={Příklad},#1}
\newtcolorbox{definiceBox}[1][]{commonlesson,title={Definice},#1}

% Spolecny box pro zadani a jeho reseni, bez mezery mezi castmi.
% Pouziti:
%   \begin{prikladReseniBox}
%     Zadani ulohy.
%     \tcblower
%     \begin{align*} ... \end{align*}
%   \end{prikladReseniBox}
% Spodni cast ma jemny sedy podklad a delici linku; lze v ni sazet
% vice radku matematiky. Obe casti maji jen jeden spolecny ramecek.
\tcbset{solutionlesson/.style={
  commonlesson,
  enhanced jigsaw,
  breakable,
  colbacklower=black!3,
  segmentation style={solid,draw=black!25,line width=0.35pt},
  fontlower=\small\color{notegray},
  before lower={\noindent\textbf{Řešení:}\enspace},
  leftlower=3mm,rightlower=3mm
}}
\newtcolorbox{prikladReseniBox}[1][]{solutionlesson,title={Příklad},#1}
\newtcolorbox{zadaniReseniBox}[1][]{solutionlesson,title={Zadání},#1}

% Pomucky pro autora pripravy: sedive poznamky; \reseni pro starsi pripravy
\newcommand{\poznamka}[1]{\par\begingroup\small\color{notegray}\itshape
  \textnormal{\textbf{Poznámka pro učitele.}}\enspace #1\par\endgroup}
\newcommand{\reseni}[1]{\par\begingroup\small\color{notegray}
  \textbf{Očekávané řešení:}\enspace #1\par\endgroup}
\newcommand{\blok}[1]{\needspace{5\baselineskip}\subsection*{#1}}
\newcommand{\vyrok}[1]{\uv{#1}}


\nastavPripravu{1. ročník}{Odmocniny}{37}{Složené výrazy a podmínky}

\begin{document}

Ve složeném výrazu nejprve vyznač základy mocnin, rozsah odmocnítek a všechny jmenovatele. Pak stanov podmínky původního zápisu. Úpravy prováděj po menších krocích a nevyčísluj velká čísla dříve, než zkrátíš společné činitele.

\begin{definiceBox}[title={Definice},breakable=false]
\textbf{Definiční obor výrazu} je množina všech hodnot proměnných, pro které jsou všechny operace v původním zápisu definovány. Výsledná úprava se s původním výrazem shoduje pouze na tomto oboru.
\end{definiceBox}


U druhé odmocniny požadujeme nezáporný radikand, u jmenovatele nenulovou hodnotu. U záporné celočíselné mocniny je základ nenulový; pro racionální mocniny používáme kladné základy podle 31. hodiny. Zjednodušení zápisu původní zákazy neruší.

\begin{prikladReseniBox}[title={Příklad 1},breakable=false]
Uprav $\frac{\sqrt{12}\sqrt{18}}{\sqrt6}$.
\tcblower
Všechny odmocniny existují a jmenovatel je kladný. Sloučením do jedné odmocniny dostáváme $\sqrt{12\cdot18/6}=\sqrt{36}=6$.
\end{prikladReseniBox}


\begin{prikladReseniBox}[title={Příklad 2},breakable=false]
Uprav $\frac{\sqrt{18a^2}}{a\sqrt2}$ pro reálné $a$.
\tcblower
Čitatel existuje pro všechna $a$, ale jmenovatel vyžaduje $a\ne0$. Částečně odmocníme a zkrátíme nenulové $\sqrt2$: $\frac{3|a|\sqrt2}{a\sqrt2}=\frac{3|a|}{a}$. Pro $a>0$ je hodnota $3$, pro $a<0$ hodnota $-3$.
\end{prikladReseniBox}


\blok{Samostatné procvičení}

\begin{zadaniReseniBox}[title={Procvičení},breakable=false]
\begin{enumerate}[label=\alph*)]
\item Uprav $\frac{(4x^{-2}y^3)^2}{8x^{-3}y}$.
\item Uprav $\sqrt{75}-\sqrt{12}+\sqrt{27}$.
\item Uprav $\frac{\sqrt{x^2}}{x}$ a uveď podmínky.
\end{enumerate}
\tcblower
\begin{enumerate}[label=\alph*)]
\item Při $x,y\ne0$ vyjde $\frac{16x^{-4}y^6}{8x^{-3}y}=\frac{2y^5}{x}$.
\item $5\sqrt3-2\sqrt3+3\sqrt3=6\sqrt3$.
\item Pro $x\ne0$ dostaneme $|x|/x$, tedy $1$ pro $x>0$ a $-1$ pro $x<0$.
\end{enumerate}
\end{zadaniReseniBox}


\poznamka{Nejprve nech celou třídu samostatně napsat podmínky P2. Rezerva propojuje více pravidel a slouží rychlejším žákům; její vynechání neomezuje základní obsah. Záměnu $\sqrt{a^2}=a$ vyvrať záporným dosazením.}


\blok{Navazující procvičování ve třídě}

\begin{zadaniReseniBox}[title={Rezerva: od jednoduché podmínky ke složené úpravě},breakable=false]
\begin{enumerate}[label=\alph*)]
\item Uprav $\frac{\sqrt{32x^2}}{2x\sqrt2}$ a uveď zvlášť výsledek pro kladné a záporné $x$.
\item Pro $a>0$ uprav $\frac{a^{3/2}\sqrt a}{a^{-1}}$.
\end{enumerate}
\tcblower
\begin{enumerate}[label=\alph*)]
\item $x\ne0$; $\frac{4|x|\sqrt2}{2x\sqrt2}=2|x|/x$, tedy $2$ pro $x>0$ a $-2$ pro $x<0$.
\item $a^{3/2+1/2+1}=a^3$.
\end{enumerate}
\end{zadaniReseniBox}

\blok{Volitelná rezerva}

\begin{prikladReseniBox}[title={Rozšíření procvičení},breakable=false]
Pro $a>0$, $b\ne0$ uprav $\left(\frac{a^{1/2}b^{-1}}{a^{-1/2}b}\right)^{-2}$.
\tcblower
Vnitřní podíl je $a^{1/2+1/2}b^{-1-1}=ab^{-2}=a/b^2$. Je nenulový, proto ho lze umocnit na $-2$. Výsledek je $b^4/a^2$ při původních podmínkách.
\end{prikladReseniBox}



\poznamka{Základ: původní příklady a samostatné procvičení, s povinným zápisem podmínek před každou úpravou. Nová dvojice je rezerva k témuž postupu; je vhodné začít prvním příkladem společně, pokud se absolutní hodnota zavádí příliš náhle. Domácí čtveřice upevňuje stejné dovednosti.}

\blok{Domácí procvičení}

\begin{zadaniReseniBox}[title={Domácí úlohy},breakable=false]
\begin{enumerate}[label=\alph*)]
\item Uprav $\frac{(3x^{-1}y^2)^2}{9x^{-3}y}$ a urči podmínky.
\item Uprav $\sqrt{80}-\sqrt{45}+\sqrt{20}$.
\item Urči podmínky a uprav $\frac{\sqrt{8a^2}}{a\sqrt2}$.
\item Pro $x>0$ uprav $(x^{1/2})^{-2}x^3$.
\end{enumerate}
\tcblower
\begin{enumerate}[label=\alph*)]
\item $xy^3$, $x,y\ne0$.
\item $4\sqrt5-3\sqrt5+2\sqrt5=3\sqrt5$.
\item $a\ne0$; $2|a|/a$, tedy $2$ pro $a>0$, $-2$ pro $a<0$.
\item $x^2$, stále $x>0$.
\end{enumerate}
\end{zadaniReseniBox}

\end{document}
